% Generated draft: requires mathematical and source review.
\documentclass[11pt]{article}
\usepackage[margin=25mm]{geometry}
\usepackage{fontspec}
\setmainfont{Latin Modern Roman}
\usepackage{amsmath,amssymb,mathtools,unicode-math}
\setmathfont{Latin Modern Math}
\usepackage{enumitem,needspace,xurl,hyperref}
\setlength{\parindent}{0pt}
\setlength{\parskip}{5pt}
\setlength{\emergencystretch}{3em}
\newcommand{\N}{\mathbb{N}}
\newcommand{\Z}{\mathbb{Z}}
\newcommand{\Q}{\mathbb{Q}}
\newcommand{\R}{\mathbb{R}}
\newcommand{\C}{\mathbb{C}}
\newcommand{\F}{\mathbb{F}}
\newcommand{\A}{\mathbb{A}}
\newcommand{\PP}{\mathbb{P}}
\newcommand{\E}{\mathbb{E}}
\newcommand{\eps}{\varepsilon}
\newcommand{\dd}{\,\mathrm{d}}
\begin{document}
\section*{UnsolvedMath ID: 700588}
\textbf{Source problem code:} catalogue entry
\subsection{Definitions and mathematical statement}
Let \(p\) be a prime and let \(G\) be a finite group of order \(p^m\), where \(m\geq 0\). Its integral group ring is \(\Z G\), with augmentation map
\[
\varepsilon:\Z G\longrightarrow\Z,\qquad \varepsilon\!\left(\sum_{g\in G}a_g g\right)=\sum_{g\in G}a_g.
\]
For the formulation below, an augmented automorphism means a unital ring automorphism \(\varphi:\Z G\to\Z G\) satisfying \(\varepsilon\circ\varphi=\varepsilon\), and a central automorphism means one that fixes the center pointwise:
\[
\varphi(z)=z\qquad\text{for every }z\in Z(\Z G),
\]
where \(Z(\Z G)\) denotes the center of the ring. The supplied record does not define these terms; in particular, the intended meaning of ``central'' requires confirmation. The question below uses the stated conventions explicitly.

Define the \(p\)-adic group ring by
\[
A=\widehat{\Z}_p G:=\varprojlim_{n\geq 1}(\Z/p^n\Z)G.
\]
A unital automorphism \(\varphi\) of \(\Z G\) fixes the integers and preserves each ideal \(p^n\Z G\), so it induces a continuous automorphism \(\widehat\varphi\) of \(A\). For \(\alpha\in\operatorname{Aut}(G)\), let \(\alpha_*\) be the induced \(\widehat{\Z}_p\)-algebra automorphism of \(A\), given by
\[
\alpha_*\!\left(\sum_{g\in G}a_g g\right)=\sum_{g\in G}a_g\alpha(g).
\]
Write \(A^\times\) for the group of units of \(A\).

\textbf{Existence question.} Does there exist a prime \(p\), a finite \(p\)-group \(G\), and a central augmented automorphism \(\varphi\) of \(\Z G\), in the senses specified above, such that
\[
\nexists\,(u,\alpha)\in A^\times\times\operatorname{Aut}(G)
\quad\text{with}\quad
\widehat\varphi(x)=\alpha_*(u x u^{-1})\quad\text{for every }x\in A?
\]
Here the displayed composition means conjugation by a unit followed by the map induced by a group automorphism, as in the supplied question. This is a statement-only formulation conditional on the terminology above; it does not resolve the source's undefined terminology.
\subsection{Short English statement}
Under the stated interpretation of ``central augmented,'' the question asks whether some finite \(p\)-group has an augmentation-preserving automorphism of its integral group ring that fixes the ring center pointwise, but whose extension to the \(p\)-adic group ring cannot be obtained by conjugating by a unit and then applying a group automorphism. The source's intended terminology requires mathematical review.
\subsection{Sources}
\begin{itemize}
\item Local UnsolvedMath registry: supplied complete source record, permanent ID \(700588\), exact source problem code \texttt{catalogue entry}, title ``Kourovka Notebook Problem 9.55.'' The registry labels the problem \texttt{open}; this is supplied metadata, not an independently verified mathematical or literature assessment.
\item The supplied attribution is to K.\ W.\ Roggenkamp, proposed in 1984, in \emph{The Kourovka Notebook (Unsolved Problems in Group Theory)}, 9th issue (1984), original problem 9.55. The record identifies the 21st issue edition (2026), arXiv:1401.0300v48, as its source edition.
\item The supplied collection description names E.\ I.\ Khukhro and V.\ D.\ Mazurov as editors of the Kourovka Notebook collection covering Issues 1--21 (1965--2026). These bibliographic details are reproduced from the record; the cited source has not been independently consulted here.
\end{itemize}
\end{document}
